%%%%%%%%%%%%%%%%%%%%%%%%%%%%%%%%%%%%%%%%%%%%%%%%%%%%%%
% Filename:      graduation_design.tex
% Author:        Junwei Wang(wakemecn@gmail.com)
% Last Modified: 2012-05-06 11:07
% Description:
%%%%%%%%%%%%%%%%%%%%%%%%%%%%%%%%%%%%%%%%%%%%%%%%%%%%%%
\documentclass[a4paper,12pt,oneside,openany]{book}
\usepackage{ctex}
\usepackage[top=1in,bottom=1in,left=1.25in,right=1.25in]{geometry}
\usepackage{graphicx}
\usepackage{ulem}
\usepackage{amsfonts,amssymb}
\usepackage{amsmath}
\usepackage{hyperref}
\usepackage[bf,small,center,indentafter,pagestyles]{titlesec}
% redefine some name
\renewcommand{\chaptername}{第\chinese{chapter}章}
\renewcommand{\bibname}{参考文献}
% redefine the chapter section
\titleclass{\chapter}{straight}
\titleformat{\chapter}[hang]
{\centering\zihao{-3}\bfseries}{第\chinese{chapter}章}{1em}{}
\titleformat{\section}
{\normalfont\zihao{4}\bfseries}{\thesection}{11pt}{\large}
% define pagehead style
\newpagestyle{firstpage}{
}
\newpagestyle{abstract}{
    \sethead{}{}{\kaishu\small\hfill 摘\qquad 要\hfill \thepage}
    \setfoot{}{}{}\headrule
}
\newpagestyle{tableofcontent}{
    \sethead{}{}{\kaishu\small\hfill 目\qquad 录 \hfill \thepage}
    \setfoot{}{}{}\headrule
}
\newpagestyle{main}{
   \sethead{}{}{\kaishu\small\hfill \chaptername\quad\chaptertitle\hfill\thepage}
   \setfoot{}{}{}\headrule
}
%%
\usepackage{titletoc}
\usepackage[titletoc]{appendix}
%定义目录格式
\titlecontents{chapter}[0em]{\normalsize}
{\thecontentslabel\hspace*{1em}}{\hspace*{-2.3em}}
{\titlerule*[0.8pc]{.}\contentspage}

\titlecontents{section}[0em]{\normalsize}
{\thecontentslabel\hspace*{1em}}{\hspace*{-2.3em}}
{\titlerule*[0.8pc]{.}\contentspage}

\titlecontents{subsection}[0em] {\normalsize}
{\thecontentslabel\hspace*{1em}}{\hspace*{-2.3em}}
{\titlerule*[0.8pc]{.}\contentspage}
%%%%%%%%%%%%%code style%%%%%%%%%%
\usepackage{xcolor}
\renewcommand{\ttdefault}{phv}
\usepackage{listings}
\renewcommand{\lstlistingname}{代码}
\lstset{
         basicstyle=\footnotesize\ttfamily, % Standardschrift
         %numbers=left,               % Ort der Zeilennummern
         numberstyle=\tiny,          % Stil der Zeilennummern
         %stepnumber=2,               % Abstand zwischen den Zeilennummern
         numbersep=5pt,              % Abstand der Nummern zum Text
         tabsize=2,                  % Groesse von Tabs
         extendedchars=true,         %
         breaklines=true,            % Zeilen werden Umgebrochen
         keywordstyle=\color{red},
            frame=b,         
         stringstyle=\color{white}\ttfamily, % Farbe der String
         showspaces=false,           % Leerzeichen anzeigen ?
         showtabs=false,             % Tabs anzeigen ?
         xleftmargin=17pt,
         framexleftmargin=17pt,
         framexrightmargin=5pt,
         framexbottommargin=4pt,
%        backgroundcolor=\color{lightgray},
         showstringspaces=false      % Leerzeichen in Strings anzeigen ?       
}
\lstloadlanguages{% Check Dokumentation for further languages ...
         %[Visual]Basic
         %Pascal
         %C
         %C++
         %XML
         %HTML
         Java
}
\usepackage{caption}
\DeclareCaptionFont{white}{\color{white}}
\DeclareCaptionFormat{listing}{\colorbox[cmyk]{0.43, 0.35, 0.35,0.01}{\parbox{\textwidth}{#1#2#3}}}
\captionsetup[lstlisting]{format=listing,justification=raggedright,labelfont=white,textfont=white, singlelinecheck=false, margin=0pt, font={sf,bf,footnotesize}}
%%%%%%%%%%%%%code style%%%%%%%%%%
% 取消中文连字符与文字间的空格
\normalspacedchars{-}
% 定义"defination"环境
\newtheorem{Definition}{定义}
\begin{document}
% titlepage
\frontmatter
\pagestyle{firstpage}
%%%%%%%%%%%%%%%%%%%%%%%%%%%%%%%%%%%%%%%%%%%%%%%%%%%%%%
% Filename:      titlepage.tex
% Author:        Junwei Wang(wakemecn@gmail.com)
% Last Modified: 2012-04-17 09:03
% Description:
%%%%%%%%%%%%%%%%%%%%%%%%%%%%%%%%%%%%%%%%%%%%%%%%%%%%%%
\vspace{30mm}
\begin{center}
% the picture
\includegraphics[scale=0.7]{sdu_pic.jpg}\\
\zihao{0} 毕业论文（设计）\\
%
\vspace{15mm}
\leftline{\heiti\zihao{-2} 设计（论文）题目：}
\heiti\zihao{-2} 密文策略--基于属性的加密\par
%
\vspace{60mm}
\zihao{4}姓\qquad 名\uline{\kaishu\makebox[50mm]{王军委}}\\
\zihao{4}学\qquad 号\uline{\kaishu\makebox[50mm]{200800300237}}\\
\zihao{4}学\qquad 院\uline{\kaishu\makebox[50mm]{软件学院}}\\
\zihao{4}专\qquad 业\uline{\kaishu\makebox[50mm]{软件工程}}\\
\zihao{4}年\qquad 级\uline{\kaishu\makebox[50mm]{2008级}}\\
\zihao{4}指导老师\uline{\kaishu\makebox[50mm]{胡程瑜}}\\
%
\vspace{5mm}
\zihao{4}\today
\end{center}

% abstract
\pagestyle{abstract}
%%%%%%%%%%%%%%%%%%%%%%%%%%%%%%%%%%%%%%%%%%%%%%%%%%%%%%
% Filename:      abstract.tex
% Author:        Junwei Wang(wakemecn@gmail.com)
% Last Modified: 2012-05-06 10:45
% Description:
%%%%%%%%%%%%%%%%%%%%%%%%%%%%%%%%%%%%%%%%%%%%%%%%%%%%%%
\centerline{\zihao{-2}\bfseries 密文策略--基于属性的加密}
\addcontentsline{toc}{chapter}{\qquad 摘\qquad 要}
\chapter*{摘\qquad 要}
在分布式系统中，只有当用户拥有某些凭据或者属性才能够访问数据。
目前，执行这种策略的唯一方式是采用一个可信的服务器来存储数据
并调整访问控制。然而，如果任一存储数据的服务器被攻破，那么数
据的保密性就会大打折扣。在本文中，我们提供一个称之为密文策略
--基于属性的加密来实现对加密数据的复杂的访问控制。通过我们的
技术，加密数据甚至可以存储在不可信的存储服务器上。此外，我们
的方法可以抵制合谋攻击。先前的基于属性的加密方案用属性来描述
被加密的数据，而将策略写入用户的的密钥里；然而在我们的系统中
属性用来描述用户的凭据，加密数据的一方来决定谁可以解密的策略
。因此，我们的加密方案在概念上更接近与传统的访问控制方法，如
基于角色的访问控制（Role-Base Access Control，RBAC）。\par
除此之外，我们提供了我们系统的实现和性能测试。\\
\textbf{关键字：} 访问控制；合谋攻击；密文策略；基于属性的加密
\addcontentsline{toc}{chapter}{\qquad ABSTRACT}
\chapter*{ABSTRACT}
In several distributed systems a user should only be %
able to access data if a user possesses a certain set of %
credentials or attributes. Currently, the only method %
for enforcing such policies is to employ a trusted server %
to store the data and mediate access control. However, %
if any server storing the data is compromised, then the %
confidentiality of the data will be compromised. In this %
paper we present a system for realizing complex access %
control on encrypted data that we call Ciphertext-Policy %
Attribute-Based Encryption. By using our techniques %
encrypted data can be kept confidential even if the %
storage server is untrusted; moreover, our methods are %
secure against collusion attacks. Previous Attribute-Based %
Encryption systems used attributes to describe %
the encrypted data and built policies into user’s keys;%
while in our system attributes are used to describe a %
user’s credentials, and a party encrypting data determines %
a policy for who can decrypt. Thus, our methods %
are conceptually closer to traditional access control %
methods such as Role-Based Access Control (RBAC).\par
In addition, we provide an implementation of our system %
and give performance measurements.\\
\textbf{Keyword:}Access Control;Collusion Attacks;%
Ciphertext-Policy;Attribute-Based Encryption

% contents
\pagestyle{tableofcontent}
\renewcommand{\contentsname}{目\quad 录}
\tableofcontents
% 绪论 
\mainmatter
\pagestyle{main}
%%%%%%%%%%%%%%%%%%%%%%%%%%%%%%%%%%%%%%%%%%%%%%%%%%%%%%
% Filename:      introduction.tex
% Author:        Junwei Wang(wakemecn@gmail.com)
% Last Modified: 2012-04-16 21:35
% Description:
%%%%%%%%%%%%%%%%%%%%%%%%%%%%%%%%%%%%%%%%%%%%%%%%%%%%%%
\chapter{绪论}
\section{基于属性加密的研究背景}
在许多情况下，当用户对敏感数据进行加密时，这时他迫切需要建立特定
的访问控制策略决定谁可以解密数据。例如，假设FBI在Knoxvile和San
Francisco的公共反腐办公室调查涉及San Francisco说客行贿Tennesseei
国会议员的指控。上级FBI人员可能想加密一份敏感的备忘录以便只有拥有
特定凭据和属性的人员可以查看。例如，上级人员可能指定下述访问结构
来访问此信息：
\centerline{\textsc{(("Public Corruption Office" \textbf{AND} 
("Knoxville" }}
\centerline{\textsc{\textbf{OR} "San Francisco"))
\textbf{OR} "Name: Charlie Eppes")}} 
通过这个访问结构，上级人员可能意味着备忘录只能
被工作在公共Knowxille或者San Francisco反腐办公室的或者名字叫做
Charlie Eppes的人员查看。\par
就像上例中所解释的，秘密数据拥有者能够根据特定的只是选择访问策略。
除此之外，这个人可能不知道其他所有能够访问被加密数据人员的精确身
份，但至少他有以属性或者凭据的方式描述这些人的方法。\par
传统方案中，这种表达式形式的访问控制必须采用可信服务器本地存储数
据的方式。服务器作为通过检查用户是否拥有某种凭据来控制访问记录或
文件的检查器被信任。然而，随着数据的增长，数据会被存储到越来越多
服务器的分布式环境中。在分布式环境中复制数据具有高性能和高可靠性
的优势。但是随着数据的增长，保证数据安全性越来越难；当数据存储在
多个服务器上，任一这些服务器被攻破的可能性大幅攀升。基于这些原因
我们需要敏感数据以加密的方式存储，以便在服务器被攻破的情况下数据
仍然保密。\par
公钥加密是一个保障存储和传输敏感数据的强有力机制。传统上，加密被
看作是用户向目标用户或者设备分享数据的一种方法。然而这种方法仅适
用于数据提供方精确的知道确切的知道要跟谁分享数据，然而更多的应用
中数据提供者想要更具接受者的凭据来制定分享机制。也就是说，大多数
现存的公钥加密方案允许一方给另一方加密数据，但是不能够高效的处理
像上述例子中更复杂的表达式形式的访问控制。\par
\section{国内外研究现状}
Sahai和Waters\cite{SW:FuzzyIBE}提出了加密算法本身决定谁可以解密被加
密的加密新思路。数据提供者用谓词$f(\cdot)$描述他想他想如何分享数
据，用户被赋予与他们的凭据$X$关联的密钥;当$f(X)=1$时，拥有凭据
$X$的用户可以解密拥有谓词$f$的密文。Sahai和Waters\cite{SW:FuzzyIBE}
提出了这种问题的正式描述，他们称之为基于属性加密\footnote{Attribute-Based 
Encryption, 简称ABE}。在ABE中，用户的平局被称为“属性”的字符
串集合表示，谓词用对这些属性的公式表示。SW\footnote{Amit Sahai和
Brent Waters}的技术受到基于身份的加密\footnote{Identity-Based 
Encryption，以下IBE}上一些工作\cite{Shamir:IBC,BF:IBE,Cocks:IBE,CHK:fspk,BB:IBE}
的启发。SW方法的一个缺点是他们最初的构造受限于处理由阈值门组成的
公式。\par
在随后的工作中，Goyal, Pandey, Sahai, and Waters\cite{GPSW:ABE}
进一步阐明了ABE的概念。尤其是，他们提出了两种互补形式的ABE。
其一是密钥策略基于属性的加密\footnote{Key-Policy ABE,简称KPABE},
属性用来描述密文，对属性的公式用来描述用户的私钥。其二是与KPABE
互补的密文策略基于属性的加密\footnote{Ciphertex-Policy ABE,简称
CPABE}，属性用来描述用户的凭据即私钥，公式则被加密放附加到密文。
除此之外。Goyal等人\cite{GPSW:ABE}提供了KPABE的一种构造。由于它
允许被加密数据的密钥能够被任意单调访问公式描述，所以该构造的代
价较高。这个系统被证明在双线型Diffie-Hellma假设下选择性安全。然
而他们把CPABE描述性的构造留做了一个开放性问题。\par
第一次明确解决CPABE问题的是Bethencourt，Sahia和Waters\cite{BSW:CPABE}。
他们描述了一个高效的可以描述的系统，这个系统允许加密者以任意单调
访问公式来表达访问谓词$f$，并取得了Gyal等人\cite{GPSW:ABE}系统的
类似表现和效率。
\section{本文的主要工作}
本文主要描述BSW的工作\cite{BSW:CPABE}。在BSW的系统中用户的私钥与
任意数目的字符串形式的属性相关联。另一方面，当一方要要加密数据时，
他们明确相关的属性上的访问结构。仅当用户的属性可以通过密文的访问
结构，他才能解密密文。从数学角度讲，访问结构被描述为单调“访问树”，
访问结构的节点由阈值门组成，叶子节点描述了属性。我们用\textbf{AND}
来构造$n\ of\ n$阈值门，用\textbf{OR}表示$1\ of\ n$阈值门。
\section{本文的主要技术}
从较高的水平上讲，本文的主要工作与最近Sahai和Waters\cite{SW:FuzzyIBE},
Goyal等人\cite{GPSW:ABE}的工作相似，但是本文却大量采用了新的技术。
KPABE，密文与描述性的属性关联，用户的私钥与策略关联（与我们的情况
恰好相反）。\textit{我们强调KPABE中，加密者对谁可以访问他所加密的
数据不具有控制权，除非她选择对数据的描述性数据。}然而，他必须保证
提供给用户适当的私钥以授权或者禁止用户解密密文。换言之，在
\cite{SW:FuzzyIBE,GPSW:ABE}中，“情报”掌握在私钥询问者手中，而不是
加密者手中。在我们的构造里，加密者决定了其他人能否访问她所加密的
数据。因此，在\cite{SW:FuzzyIBE,GPSW:ABE}中的技术不适合我们，我们
必须采用新的技术。\par
从技术的水平上讲，我们必须达到的主要目标是\textit{可抵抗合谋攻击}:
如果多个用户合谋，他们只能解密至少他们中的一个用户可以独立解密的
密文。特别提到绪论一开始提到的例子，假设工作在San Francisco反恐局
的FBI特工与他工作在New York公共反腐局的朋友合谋。我们不希望合谋者
能够解密由他们共同属性加密的备忘录。在我们的构造中，这种安全是
\textit{先决条件}。\par
在\cite{SW:FuzzyIBE,GPSW:ABE}的工作中，通过利用秘密分享方案
\cite{Shamir:SSS,Blakley:SSS}和独立地选择私密份额嵌入到每一个人的
私钥中来保证可地址合谋攻击。
\section{本文的组织结构}


% 相关工作
%%%%%%%%%%%%%%%%%%%%%%%%%%%%%%%%%%%%%%%%%%%%%%%%%%%%%%
% Filename:      related_work.tex
% Author:        Junwei Wang(wakemecn@gmail.com)
% Last Modified: 2012-04-17 10:13
% Description:
%%%%%%%%%%%%%%%%%%%%%%%%%%%%%%%%%%%%%%%%%%%%%%%%%%%%%%
%\chapter{相关工作}
%\section{IBE及其扩展}
%Sahai和Waters提出了ABE的概念的同时，提出了一个叫做模糊的基于身份的加密
%\footnote{Fuzzy-IBE,简称FIBE}的特别系统。FIBE机制建立在一些IBE思想的基础
%\cite{BF:IBE,Shamir:IBE,Cocks:IBE,CHK:fspk,BB:IBE}之上。在FIBE中，身份被
%看作是属性的集合。如果私钥拥有身份$\omega$，密文基于身份$\omega'$，当且仅
%当$\omega$和$\omega'$“设置重叠”的距离度量足够接近私钥才能解密密文。换言之，
%如果消息用属性集合$\omega'$加密，私钥由属性集合$\omega$构成，当且仅当
%$|\omega\cap\omega'|\geq d$\footnote{d在加密是已经确定}时，私钥才能解密密
%文。因此，FIBE实现了容错并适用于生物身份识别。然而，FIBE却不适用我们本文
%的主要目的即数据的访问控制。由于FIBE的主要目标是实现容错，因此它唯一支持
%的访问结构是阈值在加密时已经确定的阈值门。\par
%

% 背景 
%%%%%%%%%%%%%%%%%%%%%%%%%%%%%%%%%%%%%%%%%%%%%%%%%%%%%%
% Filename:      background.tex
% Author:        Junwei Wang(wakemecn@gmail.com)
% Last Modified: 2012-04-21 22:01
% Description:
%%%%%%%%%%%%%%%%%%%%%%%%%%%%%%%%%%%%%%%%%%%%%%%%%%%%%%
\chapter{背景}
我们首先给出CP-ABE安全性的正式定义，然后我们给出双线性映射的背景信息。如同Goyal等人
\cite{GPSW:ABE}所做，我们定义访问结构并在我们的安全性定义中使用；不同的是，在我们的
定义中属性将用来描述用户，访问结构用来标记被加密的数据。
%%%%%%%%%%%%%%%%%%%%%%%%%%%%%%%%%%%%%%%%%%%%%%%%%%%%%%%
%%%%%%%%%%%%%%%%%%%%%%%%%%%%%%%%%%%%%%%%%%%%%%%%%%%%%%%
\section{定义}
\begin{Definition}[访问结构\cite{Beimel:SSSS}]
令$\{P_1,P_2,...,P_n\}$是所有的参与方集合。集合$\mathbb{A}\subseteq 2^{\{P_1,P_2,...,P_n\}}$
是单调的是指：对于任意$B,C$，如果$B\in\mathbb{A}$且$B\subseteq C$，则$C\in\mathbb{A}$。
访问结构（特别的，单调访问结构）是$\{P_1,P_2,...,P_n\}$的非空子集，例如，
$\mathbb{A}\subseteq 2^{\{P_1,P_2,...,P_n\}}\setminus\{\emptyset\}$，
的集合$\mathbb{A}$（特别的，单调集合）。集合$\mathbb{A}$中的集合被称为授权集，不在
集合$\mathbb{A}$中的集合被称为未授权集。
\end{Definition}
\par
在我们的论述中，参与方被属性所代替，因此，访问结构$\mathbb{A}$包含了授权属性集。本文中
我们只注意单调访问结构，但是，以我们的技术实现通用的访问结构是可能的。从现在开始，除非
特别强调，访问结构是指单调访问结构。\par
%%%%%%%%%%%%%%%%%%%%%%%%%%%%%%%%
\subsection{CP-ABE方案}
一个CP-ABE机制包含四个基本的算法：Setup，Encrypt，KeyGen和Decrypt。\\
\textbf{Setup}\quad Setup算法的输入只有隐含的安全参数，输出公共参数PK和主密钥MK。\\
\textbf{Encrypt(Pk,$M$,$\mathbb{A}$)}\quad 加密算法的输入是公共参数PK，明文$M$和属性全集
上的访问结构$\mathbb{A}$。算法加密明文$M$，产生密文CT，只有用户持有满足访问结构的属性
集合才能解密消息。我们假定密文隐含访问结构$\mathbb{A}$。\\
\textbf{Key Generation(MK,$S$)}\quad
密钥生成算法的输入是主密钥MK和描述私钥的属性集合$S$，输出是私钥SK。\\
\textbf{Decrypt(PK,CT,SK)}\quad
解密算法的输入是公共参数PK，包含访问策略$\mathbb{A}$的密文CT和属性集合$S$生成的私钥SK。
如果集合$S$满足访问结构$\mathbb{A}$那么算法会解密密文CT返回消息$M$。\par
%%%%%%%%%%%%%%%%%%%%%%%%%%%%%%%%
现在我们将描述CP-ABE方案的安全模型。如同IBE方案\cite{Shamir:IBC,BF:IBE,Cocks:IBE}的安全
模型，CP-ABE方案的安全模型允许敌手询问任何比能解密密文的私钥，即在我们的安全定义中敌手
挑战访问结构$\mathbb{A}^{*}$的解密，并且可以询问不满足$\mathbb{S}^{*}$的私钥$S$。
%%%%%%%%%%%%%%%%%%%%%%%%%%%%%%%%
\subsection{CP-ABE安全模型}
现在我们给出正式的安全游戏。\\
\textbf{Setup}\quad 挑战者运行Setup算法并将得到的公共参数PK交给敌手。\\
\textbf{阶段1}\quad 敌手反复生成对应于属性集$S_1,...,S_{q_1}$的私钥。\\
\textbf{Challenge}\quad 敌手提供两个等长的信息$M_0$和$M_1$。除此之外，敌手提供一个阶段1
中所有属性集合$S_1,...,S_{q1}$都不满足的挑战访问结构$\mathbb{A^*}$。挑战者随机选取
$b\in{0,1}$，在访问结构$\mathbb{A^*}$下加密$M_b$，并将密文$CT^*$敌手。\\
\textbf{阶段2}\quad 在属性集合$S_{q_1},...,S_q$不能满足对应的挑战访问结构的限制下重复
阶段1。\\
\textbf{Guess}\quad 敌手输出对$b$的猜测$b'$。\par
%%%%%%%%%%%%%%%%%%%%%%%%%%%%%%%%
\vspace{3mm}
在上述游戏中敌手$\mathcal{A}$的优势被定义为$\mathrm{Pr}[b'=b]-\frac{1}{2}$。我们注意到这个模型可以通过在阶段1和阶段2中允许解密敌手询问的明文来扩展为可以处理选择明文攻击的情形。
\begin{Definition}
在上述游戏中，如果多项式时间内敌手最多有可以忽略的优势，那么CP-ABE方案是安全的。
\end{Definition}
%%%%%%%%%%%%%%%%%%%%%%%%%%%%%%%%%%%%%%%%%%%%%%%%%%%%%%%
%%%%%%%%%%%%%%%%%%%%%%%%%%%%%%%%%%%%%%%%%%%%%%%%%%%%%%%
\section{双线性映射}
我们提出一些与群相关的高效可计算双线型映射的事实。\par
设$\mathbb{G}_0$和$\mathbb{G}_1$是阶为素数$p$的两个乘法循环群，$g$是$\mathbb{G}_0$
的生成元，$e:\mathbb{G}_0\times \mathbb{G}_0\to \mathbb{G}_1$。双线性映射$e$具有以
下性质：
\begin{enumerate}
\setlength{\itemsep}{0pt}
\item 双线性：对于任意$u,v\in \mathbb{G}_0$和任意$a,b\in \mathbb{Z}_p$，我们有
$e(u^a,v^b) =e(u,v)^{ab}$
\item 非退化性：$e(g,g)\neq 1$
\end{enumerate}
\par
如果$\mathbb{G}_0$中的群操作和双线性映射
$e:\mathbb{G}_0\times \mathbb{G}_0\to\mathbb{G}_1$都是高效可计算的，我们称
$\mathbb{G}_0$为双线性群。请注意由于$e(g^a,g^b)=e(g,g)^{ab}=e(g^b,g^a)$，
双线性映射$e$具有对称性。

% 构造
%%%%%%%%%%%%%%%%%%%%%%%%%%%%%%%%%%%%%%%%%%%%%%%%%%%%%%
% Filename:      construction.tex
% Author:        Junwei Wang(wakemecn@gmail.com)
% Last Modified: 2012-05-03 10:28
% Description:
%%%%%%%%%%%%%%%%%%%%%%%%%%%%%%%%%%%%%%%%%%%%%%%%%%%%%%
\chapter{CP-ABE构造}
本章中我们提供了BSW方案的构造。我们首先介绍描述密文的访问树和描述私钥的属性的模型，
其次我们给出我们方案的描述，最后我们给出安全性，高效性的讨论。
%%%%%%%%%%%%%%%%%%%%%%%%%%%%%%%%%%%%%%%%%%%%%%%%%%%%%%
%%%%%%%%%%%%%%%%%%%%%%%%%%%%%%%%%%%%%%%%%%%%%%%%%%%%%%
\section{模型}
在我们的构造中私钥被描述性的属性集合$S$所确定。一方想解密数据必须拥有满足访问树结构的
私钥。\par
访问树中的每个内部节点是一个阈值门，叶子节点与属性关联。（请注意，这样的描述是非常
具有表达力的。例如，我们分别用2 of 2和1 of 2的阈值门来表达树中“AND”和“OR”门。）当且
仅当私钥的属性分配到树的节点并使树能满足，持有私钥的用户将能够解密密文。尽管在我们的
方案里，属性用来标记私钥，我们还是用与\cite{GPSW:ABE}相同的概念来描述访问树。\par
\vspace{5mm}
\noindent\textbf{访问树$\mathcal{T}$}\quad 用$\mathcal{T}$代表访问结构，用通过他的孩
子和阈值描述阈值门来代表树的每个非叶子节点。设节点$x$的孩子数目为$num_x$，阈值为$k_x$，
那么$0<k_x\leq num_x$。当$k_x=1$，阈值门是或门；当$k_x=num_x$，阈值门是与门。树的每个
叶子节点$x$用一个属性和阈值$k_x=1$描述。\par
为了方便的使用访问树，我们定义了一些函数。函数$\textbf{parent}(x)$表示节点$x$的
父节点,仅当$x$是叶子节点时，函数$\textbf{att}(x)$被定义为$x$所描述的属性。访问树
$\mathcal{T}$同时定义了每个节点的子节点的被标记为从1到$num$的顺序索引值，函数
$\textbf{index}(x)$定义为与节点$x$关联的这样一个索引值，且对一个给定的密钥,索引值
以任意的方式唯一的赋值给访问结构中的节点。\par
\vspace{5mm}
\noindent\textbf{满足访问树}\quad 令$\mathcal{T}$是根为$r$的访问树，$\mathcal{T}_x$
表示以$x$为根的$\mathcal{T}$的子树，因此$\mathcal{T}$等同于$\mathcal{T}_r$。如果属性
集合$\gamma$满足访问树$\mathcal{T}_x$，我们表示为$\mathcal{T}_x(\gamma)=1$。我们用如
下递归的方式计算$\mathcal{T}_x(\gamma)$：如果$x$是非叶子节点，对于$x$的所有孩子，
计算$\mathcal{T}_{x'}(r)$，当且仅当至少$k_x$个孩子返回1,$\mathcal{T}_x(\gamma)$返回1；
如果$x$是叶子节点，那么当且仅当$\textbf{attr}(x)\in \gamma$时，$\mathcal{T}_x(\gamma)$
返回1。
%%%%%%%%%%%%%%%%%%%%%%%%%%%%%%%%%%%%%%%%%%%%%%%%%%%%%%
%%%%%%%%%%%%%%%%%%%%%%%%%%%%%%%%%%%%%%%%%%%%%%%%%%%%%%
\section{构造}
设$\mathbb{G}_0$是阶为素数$p$的双线性映射，$g$是$\mathbb{G}_0$的生成元，
$e:\mathbb{G}_0\times \mathbb{G}_0\to\mathbb{G}_1$是双线性映射。安全参数$\kappa$决定
群的规模。同时，我们定义了拉格朗日系数$\Delta_{i,S}$，$i\in\mathbb{Z}_p$，$S$是集合
$\mathbb{Z}_p$中的元素：$\Delta_{i,S}(x)=\prod_{j\in S,j\neq i}\frac{x-j}{i-j}$。
哈希函数$H:{0,1}^*\to \mathbb{G}_0$是random oracle，这个函数可以把任意用二进制串
描述的的属性映射到随机的群元素。我们的构造如下。\par
%%%%%%%%%%%%%%%%%%%%%%%%%%%%
\vspace{5mm}
\noindent\textbf{Setup}\quad 
Setup算法选择生成元为$g$，阶为素数$p$的双线性映射
$\mathbb{G}_0$，指数$\alpha,\beta\in\mathbb{Z}_p$。生成公钥
$$\textrm{PK} = \mathbb{G}_0,g,h=g^{\beta},f=g^{1/\beta},e(g,g)^{\alpha}$$
和主密钥$\textrm{MK} = \beta,g^\alpha$。（注意$f$仅用于委托。）\par
%%%%%%%%%%%%%%%%%%%%%%%%%%%%
\vspace{5mm}
\noindent\textbf{Encryption(PK,$M$,$\mathcal{T}$)}\quad
加密算法在树访问结构$\mathcal{T}$下加密消息$M$。算法首先为$\mathcal{T}$每个节点$x$
（包括叶子节点）选择一个多项式$q_x$。从树的根节点$R$开始，自上而下选择多项式。节点
$x$的多项式$q_x$的度$d_x$比该节点的阈值$k_x$少一，即$d_x=k_x-1$。\par
算法从根节点$R$开始选择随机数$s\in\mathbb{Z}_p$，并设置$q_R(0)=s$。然后，算法随机
选择多项式$q_R$上的$d_R$个点来完全定义$q_R$。对于其他的顶点$x$，令
$q_x(0)=q_{parent(x)}(index(x))$，随机选择其它$d_x$个顶点来完全定义$q_x$。\par
设$\mathcal{T}$中所有叶子节点的集合为$Y$，那么在给定的树形访问结构$\mathcal{T}$下计算
密文
\begin{equation*}
\begin{split}
\textrm{CT}=&(\mathcal{T},\widetilde{C}=Me(g,g)^{\alpha s},C=h^s,\\
&\forall y\in Y:\quad C_y=g^{q_y(0)},C_y'=H(\textbf{att}(y))^{q_y(0)}).
\end{split}
\end{equation*}
%%%%%%%%%%%%%%%%%%%%%%%%%%%%
\vspace{5mm}
\noindent\textbf{KeyGen(MK,$S$)}\quad
密钥生成算法的输入是属性集合$S$，输出为被$S$所标记的密钥。算法首先选择随机数
$r\in \mathcal{Z}_p$，然后对每一个$j\in S$选择随机数$r_j\in \mathcal{Z}_p$。最后计算出
私钥
\begin{equation*}
\begin{split}
\textrm{SK}=&(D=g^{(\alpha+\gamma)/\beta},\\
&\forall j\in S:\quad D_j=g^r\cdot H(k)^{rj},D_j'=g^{rj}).
\end{split}
\end{equation*}
%%%%%%%%%%%%%%%%%%%%%%%%%%%%
\vspace{5mm}
\noindent\textbf{Decrypt(CT,SK)}\quad
我们的解密算法是一个递归算法。为了便于讨论，我们提出了解密算法的最简单形式，并在下一小
节中提出了潜在的性能改进。\par
我们首先定义递归算法Decrypt(CT,SK,$x$)，它用密文
$\textrm{CT}=(\mathcal{T},\widetilde{C},C,\forall y\in Y:\quad C_y,C_y')$，与属性集合
$S$关联的私钥SK，$\mathcal{T}$中的节点$x$作为输入。\par
当节点$x$是叶子节点，令$i=\textbf{att}(x)$，如果$i\in S$，那么
\begin{equation*}
\begin{split}
DecryptNode(CT,SK,x) &= \frac{e(D_i,C_x)}{e(D_i',C_x')}\\
&=\frac{e(g^r\cdot H(i)^{r_i},g^{q_x(0)})}{e(g^{r_i},H(i)^{q_x(0)})}\\
&=e(g,g)^{rq_x(0)}.
\end{split}
\end{equation*}
如果$i\notin S$，那么DecryptNode(CT,SK,$x$)$=\perp$。\par
现在我们考虑$x$是非叶子节点时的递归情况。算法DecryptNode(CT,SK,$x$)工作方式如下：
对于$x$的所有字节点$z$，计算$F_z=\textrm{DecryptNode(CT,SK,}z\textrm{)}$。
令$S_x$为$k_x$大小的满足$F_z\neq \perp$的子节点$z$的集合。如果不存在这样的集合，
那么这个节点不满足，且函数返回$\perp$。\par
否则，我们计算
\begin{equation*}
\begin{split}
F_x &=\prod_{z\in S_x}F_x^{\Delta_{i,S_x'(0)}}
\quad \text{,其中,$i=index(z),S_x'=\{index(z),z\in S_x\}$}\\
&=\prod_{z\in S_x}(e(g,g)^{r\cdot q_z(0)})^{\Delta_{i,S_x'(0)}}\\
&=\prod_{z\in S_x}(e(g,g)^{r\cdot q_{parent(z)}(index(z))})^{\Delta_{i,S_x'(0)}}\\
&=\prod_{z\in S_x}e(g,g)^{r\cdot q_z(0)\cdot \Delta_{i,S_x'(0)}}\\
&=e(g,g)^{r\cdot q_x(0)} \qquad \text{(使用多项式插值)}
\end{split}
\end{equation*}
并返回结果。\par
在定义了DecryptNode函数之后，我们定义解密算法。算法首先调用Decrypt(CT,SK,$R$)，
$R$是树$\mathcal{T}$的根节点。如果树满足$S$，我们令
$$A=\textrm{DecryptNode(CT,SK,}R\textrm{)}=e(g,g)^{rq_R(0)}=e(g,g)^{rs}$$.
现在算法通过下面计算解密
$$\widetilde{C}/\left(\frac{e(C,D)}{A}\right)=
\widetilde{C}/\left(\frac{e(h^s,g^{(\alpha+\gamma)/\beta})}{e(g,g)^{rs}}\right)=M$$
%%%%%%%%%%%%%%%%%%%%%%%%%%%%%%%%%%%%%%%%%%%%%%%%%%%%%%
%%%%%%%%%%%%%%%%%%%%%%%%%%%%%%%%%%%%%%%%%%%%%%%%%%%%%%
\section{直觉安全和效率}
现在我们将提出我们方案的直觉安全性和我们方案的效率的讨论。
\subsection{直觉安全}
如同之前的ABE方案，在设计我们的方案时面临的最大挑战是如何抵制合谋攻击。我们采用了
Sahai和Waters方案\cite{SW:FuzzyIBE}中的通过随机化用户私钥的方案来组织合谋；所不同的是，
我们把秘密分享的思想嵌入到密文而不是私钥中。攻击者必须恢复$e(g,g)^{\alpha s}$才能够
解密。为了做到这一点，攻击者必须从用户的私钥中配对$C$，$D$两部分。这回得到想要的值
$e(g,g)^{\alpha s}$，但是被某些值$e(g,g)^{rs}$所蒙蔽。当且进当某用户拥有密钥的部分满足
嵌入到密文中的秘密分享，这个值才不会被蒙蔽。合谋攻击不会在上述过程中起到任何作用，因为
用户私钥的随机性随机化了盲值。\par
尽管我们的方案在选择明文攻击情况下是安全的，我们方案可以高效的扩展为选择密文攻击。
这个过程需要想Fujisaki-Okamoto变换\cite{FO}一样采用random oracle工具。
\subsection{效率}
密钥生成算法和加密算法的效率都非常容易计算。加密算法对密文访问树每个叶子节点需要两次
幂运算，密文的规模中包括每个树叶子节点的两个群元素。密钥生成算法对用户的每个属性需要
两次指数运算，私钥包含对每个属性的两个群元素。最简单的形式，解密算法需要对与私钥属性
匹配的访问树叶子节点两次配对，对沿着上述叶子到根节点的路径节点的
（至多\footnote{如果路径中存在不满足的内部节点，幂运算会减少})一次指数元算。然而，
可能存在着不止一种满足策略的方式，一个更加智能的算法可能会沿着这条线路优化，我们
会在第四章中介绍多种提升性能的方法。

% 实现
%%%%%%%%%%%%%%%%%%%%%%%%%%%%%%%%%%%%%%%%%%%%%%%%%%%%%%
% Filename:      implementation.tex
% Author:        Junwei Wang(wakemecn@gmail.com)
% Last Modified: 2012-05-02 18:11
% Description:
%%%%%%%%%%%%%%%%%%%%%%%%%%%%%%%%%%%%%%%%%%%%%%%%%%%%%%
\chapter{实现}
在本章中我们将要讨论实现第三章构造的一些实际问题，包括一些优化，我们所开发的工具包
的描述以及性能的测量。
\section{解密算法的效率改进}
尽管没有必要通过新的技术减少setup、密钥生成和加密算法的群运算，但是这却可以大幅提升
解密算法的性能。我们将在次介绍这些改进，并在4.3节中给出他们的作用。\par
\vspace{5mm}
\textbf{优化解密策略}\quad 第三章出的递归算法中，被对私钥属性匹配的叶子节点进行两次
配对，对沿着节点到根的每个内部节点（不包括根节点）至多一次指数运算。递归部分的最后一
步会有一次额外的配对。当然，每个阈值为$k$的内部节点，只保留他的$k$个孩子。考虑到在
提前叶节点满足和选择可以满足整个访问树的自己的时间，我们可能避免估计DecryptNode，因为
结果最终不被使用。\par
更精确的讲，设$M$是访问树$\mathcal{T}$的节点的子集。我们定义函数$\mathbf{restrict}(
\mathcal{T},M)$是从$\mathcal{T}$删除以下节点（同时保持阈值不变）形成的访问树。
首先，我们删除所有不在$M$中的节点，下一步我们删除沿着孩子数小于阈值$k_x$的内部节点
$x$一直到$\mathcal{T}$的根节点的所有节点。重复上述步骤直至没有节点可以删除，这时侯
得到的就是$\mathbf{restrict}(\mathcal{T},M)$。因此给定访问树$\mathcal{T}$和满足访问树
的属性集$\gamma$，自然的问题是选择集合$M$使得$\gamma$满足$\mathbf{restrict}(
\mathcal{T},M)$并且$M$中的叶子数目是最少的（考虑到配对操作的代价是最昂贵的）。这可以通过对树的一次遍历的递归算法很容易的实现。算法DecryptNode在
$\mathbf{restrict}(\mathcal{T},M)$上会取得与原来相同的结构。\par
\vspace{5mm}
\textbf{直接计算}\quad
进一步性能改进，可以放弃DecryptNode函数并采取更直接的计算。直观的讲，我们想象展平
树上对DecrypNode的递归调用，然后合并指数运算到每一个（用过的）叶子节点上。更精确的讲，
令$\mathcal{T}$是根为$r$的访问树，$\gamma$是属性集合，$M \subseteq \mathcal{T}$使得
$\gamma$满足$\mathbf{restrict}(\mathcal{T},M)$。同时假设$M$是最小的，即没有内部节点的
孩子数比阈值大。设$L \subseteq M$是$M$中的叶子节点。那么对于每一个$l \in L$，
$l$到$r$的路径表示为
$$\rho(l)=(l,\mathbf{parent}(l),\mathbf{parent}(\mathbf{parent}(l)),...r) .$$
同时，节点$x$的兄弟节点（包括$x$）表示为
$$\mathbf{sibs}(x)=\{y\mid \mathbf{parent}(x)=\mathbf{parent}(y)\}\quad.$$
在上述概念的基础上，我们可以继续直接计算Decrypt(CT,SK,$r$)的结果。首先，对每个
$l\in L$计算
\begin{align*}
z_l = \prod_{x\in \rho(l)\atop x\neq r}\quad \text{其中$i=\mathbf{index}(x),
S=\{\mathbf{y}|y\in\mathbf{sibs}(x)\}$ ,}
\end{align*}
然后计算
$$\mathrm{DecryptNode(CT,SK,}r\mathrm{)}
=\prod_{l\in L\atop i=\mathbf{att}(l)}
(\frac{e(D_i,C_l)}{e(D_i',C_l')})^{z_l}$$




% 附录
%%%%%%%%%%%%%%%%%%%%%%%%%%%%%%%%%%%%%%%%%%%%%%%%%%%%%%
% Filename:      appendix.tex
% Author:        Junwei Wang(wakemecn@gmail.com)
% Last Modified: 2012-05-05 10:47
% Description:
%%%%%%%%%%%%%%%%%%%%%%%%%%%%%%%%%%%%%%%%%%%%%%%%%%%%%%
\titleformat{\chapter}[display]{\centering\normalfont\Large\bfseries}
{附录~\Alph{chapter}}{11pt}{\Large}
\renewcommand{\thechapter}{附录\Alph{chapter}.}
\begin{appendices}
\chapter{CPABE使用演示}
\lstinputlisting[caption=DemoForCpabe.java]{DemoForCpabe.java}
\end{appendices}


\addcontentsline{toc}{chapter}{\qquad 参考文献}
\bibliographystyle{plain}
\bibliography{graduation_design}
\end{document}
